\section{Examples}\label{sec:examples}

We now work out the winners for a few switching graphs, and the second-order
constants where they are explicit. In every example $Q$ is symmetric. So the
admissible faces are the connected faces with positive start mass
(Lemma~\ref{lem:admissible}), Lemma~\ref{lem:symmetric} applies, and
Theorem~\ref{thm:window} applies at every tie. We call the faces in
$\mathcal F_\tau$ the \emph{winners} at $\tau$, and their list as $\tau$ runs
from $0$ to $\infty$ the \emph{first-order winner sequence}. By
Theorem~\ref{thm:first}(iii), for fixed $\tau$ and large $N$ every global mode
lies on a winner.

\subsection{The complete graph and the start}\label{sec:ex-complete}

\begin{proposition}[complete graph and single jumps]\label{prop:complete}
\begin{enumerate}[(a)]
\item Let $Q=J-dI$, the complete graph $K_d$ with unit rates. Then
$\lambda_F=d-\lvert F\rvert$ for every face $F$, so
$\Phi_\tau(F)=d\tau-1+\lvert F\rvert(1-\tau)$. At $\tau=1$ every face has
$\Phi_1(F)=d-1$. For $\tau<1$ the winners are the admissible vertices, and for
$\tau>1$ the only winner is $[d]$.
\item Let $Q$ be symmetric and irreducible, and let $x\in[d]$. Then
$\lambda_{[d]\setminus\{x\}}\le q_x/(d-1)$, with equality if and only if
$q_{yx}=q_x/(d-1)$ for every $y\ne x$.
\item Let $Q$ be symmetric and irreducible and let the start have full support.
Suppose that the first-order winner sequence is a single jump, that is, for some
$\tau_1>0$ the winners are vertices for $0<\tau<\tau_1$ and $[d]$ alone for
$\tau>\tau_1$. Then every vertex $x$ that wins for small $\tau$ has
$q_{yx}=q_x/(d-1)$ for every $y\ne x$.
\item Let $Q$ have unit rates on a connected graph $G_Q$ and let the start have
full support. Then the first-order winner sequence is a single jump if and only
if $G_Q$ is complete.
\end{enumerate}
\end{proposition}

\begin{proof}
(a) On a face with $k$ states, $-Q_F=dI-J$ has the eigenvalue $d-k$ on
$\mathbf 1_F$ and $d$ on the vectors with zero sum. So
$\Phi_\tau(F)=\tau(d-k)+k-1$. It equals $d-1$ at $\tau=1$, increases with $k$
for $\tau<1$ and decreases with $k$ for $\tau>1$.

(b) Put $R=[d]\setminus\{x\}$. In Lemma~\ref{lem:symmetric}(a) the test vector
$\mathbf 1_R/\sqrt{d-1}$ has $\mathcal E=\sum_{y\in R}q_{yx}/(d-1)=q_x/(d-1)$,
by symmetry. This gives the bound. Note that $(-Q_R\mathbf 1_R)_y=q_{yx}$ for
$y\in R$. If equality holds, then $\mathbf 1_R$ minimizes the Rayleigh quotient,
so it is an eigenvector of $-Q_R$ for $\lambda_R$, and every $q_{yx}$ equals
$\lambda_R$. Conversely, let $q_{yx}=c$ for all $y\in R$. Since
$\mathcal E(\lvert\varphi\rvert)\le\mathcal E(\varphi)$, there is an
eigenvector $\varphi\ge0$ of $-Q_R$ for $\lambda_R$. Then
$c\langle\varphi,\mathbf 1_R\rangle=\langle\varphi,-Q_R\mathbf 1_R\rangle
=\lambda_R\langle\varphi,\mathbf 1_R\rangle$, so $c=\lambda_R$.

(c) A vertex $x$ that wins for small $\tau$ has the least exit rate, so it wins
on all of $(0,\tau_1)$. The least exponent $\min_F\Phi_\tau(F)$ is continuous
in $\tau$, so
$\tau_1q_x=d-1$. Let $C$ be a component of $[d]\setminus\{x\}$ with
$\lambda_C=\lambda_{[d]\setminus\{x\}}$. It is admissible, so
$\tau_1\lambda_C+\lvert C\rvert-1\ge d-1$, and $\lvert C\rvert\le d-1$ gives
$\lambda_C\ge1/\tau_1=q_x/(d-1)$. So equality holds in (b).

(d) If $G_Q$ is complete, then (a) gives a single jump. Conversely, with unit
rates $q_x$ is the degree of $x$, and by (c) a vertex of least degree is joined
to every other state. So $G_Q$ is complete.
\end{proof}

Part (b) is the reason for (d). For a vertex $x$ of least exit rate, the point
of $[d]\setminus\{x\}$ lies on or to the left of the chord from $\{x\}$ to
$[d]$, and on it only when every other state feeds $x$ at the same rate.
Otherwise a face with between $2$ and $d-1$ states wins on an interval.
Proposition~\ref{prop:directjump}(a) is the analogue for sticky kernels. On
$K_d$ itself the first order decides nothing at $\tau=1$. With a uniform start
the second order makes the mode jump from a vertex straight to $[d]$
\cite[Cor.~3.7]{paperB}, as Corollary~\ref{cor:paperB} explains, although at
finite $N$ an intermediate face can hold the mode \cite[Ex.~3.8]{paperB}. The
start enters $K_F$ only through the factor $(\mu\cdot r_F)(l_F\cdot\mathbf 1)$,
and on $K_3$ this can open a window for an edge. The reason is that a start with
no mass on state $3$ raises the constants of the vertices $1,2$ and of the edge
$\{1,2\}$ by the same factor, compared with the uniform start, and leaves the
full face alone. The edge line has the middle slope, so it can rise above the
point where the other 2 lines cross.

\begin{proposition}[the start matters on $K_3$]\label{prop:K3start}
Let $Q=J-3I$, let $\mu$ be any start, and put $T=L-\frac12\log L+t$. At
$\tau^*=1$ all admissible faces tie, and the lines of Theorem~\ref{thm:window}
for the admissible faces are
\[
\ell_{\{i\}}(t)=\log\mu_i-2t,\qquad
\ell_{\{i,j\}}(t)=\log\Bigl((\mu_i+\mu_j)\sqrt{2/\pi}\Bigr)-t,\qquad
\ell_{[3]}(t)=\log\frac{3^{5/2}}{4\pi}.
\]
\begin{enumerate}[(a)]
\item There is a nonempty open interval of $t$ on which $\max_F\ell_F(t)$ is
attained only by edges if and only if
$\max_{i<j}(\mu_i+\mu_j)^2>\frac{3^{5/2}}8\max_l\mu_l$. The edges that attain it
are then those with the largest $\mu_i+\mu_j$.
\item For $\mu=(\frac12,\frac12,0)$ this interval is
\[
(t_{ve},t_{ef})=\Bigl(\tfrac12\log\tfrac\pi8,\
\tfrac12\log\tfrac2\pi-\log\tfrac{3^{5/2}}{4\pi}\Bigr)=(-0.4673558,\,-0.4412978),
\]
of width $0.0260580$. For the uniform start there is no such interval.
\item Let $\mu=(\frac12,\frac12,0)$. For every compact
$J\subset(t_{ve},t_{ef})$ and all large $N$, every global mode with $t\in J$ has
support $\{1,2\}$. For every compact $J'$ disjoint from $[t_{ve},t_{ef}]$ and
all large $N$, no global mode with $t\in J'$ has support $\{1,2\}$.
\end{enumerate}
\end{proposition}

\begin{proof}
By Proposition~\ref{prop:complete}(a) all admissible faces tie at $\tau^*=1$,
and Theorem~\ref{thm:window} applies since $K_3$ has symmetric support. Each
$Q_F$ has constant row sums, so $r_F=\mathbf 1_F$, $l_F=\mathbf 1_F/\lvert
F\rvert$ and $(\mu\cdot r_F)(l_F\cdot\mathbf 1)=\mu(F)$. This is the only
factor of $K_F$ that depends on the start, so Corollary~\ref{cor:paperB} gives
$K_F=\mu(F)\lvert F\rvert^{\lvert F\rvert-1/2}(4\pi)^{-(\lvert F\rvert-1)/2}$,
and $\lambda_F=3-\lvert F\rvert$ gives the lines.

(a) Let $S=\max_{i<j}(\mu_i+\mu_j)$ and $\bar\mu=\max_l\mu_l$. Only the top line
of each slope matters, namely $v(t)=\log\bar\mu-2t$,
$e(t)=\log(S\sqrt{2/\pi})-t$ and $f=\log(3^{5/2}/(4\pi))$. Let $v(t_{vf})=f$.
If $e(t_{vf})>f$, then $e$ is the strict maximum near $t_{vf}$. Otherwise $e\le
v$ for $t\le t_{vf}$ and $e\le f$ for $t\ge t_{vf}$, because $e-v$ increases
and $e-f$ decreases. As $e^{t_{vf}}=(\bar\mu e^{-f})^{1/2}$, the condition
$e(t_{vf})>f$ reads $S^2\cdot\frac2\pi>\bar\mu e^f$, which is the criterion.

(b) Here $S=1$, $\bar\mu=\frac12$ and $1>3^{5/2}/16=0.9743$. The ends solve $e=v$
and $e=f$. For the uniform start $S^2=\frac49<\frac{3^{5/2}}8\cdot\frac13=0.6495$.

(c) By Theorem~\ref{thm:window}(i), $\log M_F=-2L+\log L+\ell_F(t)+O(\log L/L)$
for every admissible face, uniformly for $t$ in compacts. On $J$ the line of
$\{1,2\}$ beats every other line by a fixed margin, and on $J'$ some other line
beats it by a fixed margin.
\end{proof}

The window is thin, since the criterion holds with ratio $0.97428$. At finite
$N$ its ends also drift by $O(\log L/L)$ (Figure~\ref{fig:window}). At the
left end the balanced composition of $\{1,2\}$ overtakes a vertex. By
Proposition~\ref{prop:pair} this is exactly the binary crossing of
\cite[Thm.~5.1]{paperA} with $u=h/(1-2h)$. The equation of
Theorem~\ref{thm:crossing} with the term of Proposition~\ref{prop:thirdorder},
here $T+\frac12\log T=L+\frac12\log\frac\pi8+\frac1{8T}$, gives this left end to
within $2\times10^{-5}$ at $N=10^{20}$.

\begin{figure}[t]
\centering
\includegraphics[width=\textwidth]{fig_window}
\caption{(a) The lines $\ell_F(t)$ of Proposition~\ref{prop:K3start} on $K_3$
with the start $(\frac12,\frac12,0)$. The vertex $\{3\}$ has no start mass and
is not admissible. The upper envelope is bold, and the edge window
$(-0.4674,-0.4413)$ is shaded. (b) The ends $t_{ve}$ and $t_{ef}$ of the window
at $N=10^6,10^8,\dots,10^{20}$, computed from the exact face maxima, against
$(\log L)/L$. The stars are the limits from (a). The dotted curve solves
$T+\frac12\log T=L+\frac12\log\frac\pi8+\frac1{8T}$ for the left end. At
$N=10^{20}$ the window is $(-0.43825,-0.41345)$.}
\label{fig:window}
\end{figure}

\subsection{Cycles and paths}\label{sec:ex-cycles}

Let $C_d$ ($d\ge3$) be the cycle with $q_{i,i\pm1}=1$, indices mod $d$. An
\emph{arc} is a set of $k\le d-1$ consecutive states. For $k\ge1$ put
\[
\lambda_k=2-2\cos\vartheta_k,\quad \vartheta_k=\frac\pi{k+1},\qquad
\tau_k=\frac1{\lambda_{k-1}-\lambda_k}\ (k\ge2),\qquad D_k=k+\lambda_k\tau_k .
\]

\begin{lemma}[arcs]\label{lem:arcs}
\begin{enumerate}[(a)]
\item On $C_d$ an arc of $k$ states has exit rate $\lambda_k$. A proper face that
is not an arc is not admissible, and its exponent is larger than that of its
longest arc.
\item The map $k\mapsto\lambda_k$ is strictly decreasing and strictly convex on
$k\ge1$. Hence $\tau_2<\tau_3<\cdots$.
\item $D_2=3$, and $\frac32k-\frac12<D_k<\frac32k$ for every $k\ge3$. In
particular $D_k$ is not an integer for $k\ge3$.
\end{enumerate}
\end{lemma}

\begin{proof}
(a) For $k\le d-1$ the graph on an arc is the path $P_k$, and
$-Q_F=2I-\operatorname{Adj}(P_k)$ has the eigenvalues $2-2\cos(j\pi/(k+1))$,
$1\le j\le k$. A proper face that is not an arc has at least $2$ components,
each an arc. So it is not admissible (Lemma~\ref{lem:admissible}), and by (b)
its exit rate is that of its longest component, which has fewer states.

(b) With $g(x)=2-2\cos(\pi/x)$ we have $\lambda_k=g(k+1)$, and for $x\ge2$
\[
g'(x)=-\frac{2\pi\sin(\pi/x)}{x^2}<0,\qquad
g''(x)=2\pi\Bigl(\frac{\pi\cos(\pi/x)}{x^4}+\frac{2\sin(\pi/x)}{x^3}\Bigr)>0.
\]

(c) Put $a=\pi/(2k)$, $b=\pi/(2(k+1))$ and
$\rho=\lambda_k/\lambda_{k-1}=(\sin b/\sin a)^2$. Since
$\lambda_k\tau_k=\rho/(1-\rho)$, the claim is
$\frac{k-1}{k+1}<\rho<\frac k{k+2}$. For $k=2$,
$\rho=\sin^2(\pi/6)/\sin^2(\pi/4)=\frac12$ gives $D_2=3$. As $\sin x/x$
decreases on $(0,\pi]$, $\sin b/\sin a\ge b/a=k/(k+1)$, so
$\rho\ge k^2/(k+1)^2>(k-1)/(k+1)$. For the upper bound let $k\ge3$, $m=k+1$ and
$\operatorname{sinc}x=\sin x/x$. The claim reads $\operatorname{sinc}
b/\operatorname{sinc}a<m/\sqrt{m^2-1}$, and the right side exceeds
$1+1/(2m^2)$. By $1-x^2/6\le\operatorname{sinc}x\le1-x^2/6+x^4/120$ it
suffices that $(1-\frac{b^2}6+\frac{b^4}{120})-(1-\frac{a^2}6)\le
(1-\frac{a^2}6)/(2m^2)$, that is
\[
\frac{\pi^2}{12}\cdot\frac{2m-1}{(m-1)^2}+\frac{\pi^4}{960\,m^2}\le
1-\frac{\pi^2}{24(m-1)^2}.
\]
The left side decreases in $m$ and the right side increases. At $m=4$ they are
$0.6460$ and $0.9543$. Finally, $(\frac32k-\frac12,\frac32k)$ has length
$\frac12$ and an integer end, so it contains no integer.
\end{proof}

\begin{theorem}[cycle cascade]\label{thm:cycle}
Let the start have full support.
\begin{enumerate}[(a)]
\item On $C_3=K_3$ the winners are the vertices for $\tau<1$ and $[3]$ for
$\tau>1$, and all 3 sizes tie at $\tau=1$.
\item Let $d\ge4$ and $K=\lfloor2d/3\rfloor$. The winners are the $d$ vertices
for $0<\tau<1$, the $d$ arcs of $k$ states for $\tau_k<\tau<\tau_{k+1}$
($2\le k\le K-1$), the $d$ arcs of $K$ states for $\tau_K<\tau<\tau_{\rm
full}$, and $C_d$ alone for $\tau>\tau_{\rm full}=\lceil d/3\rceil/\lambda_K$.
All these intervals are nonempty. Arcs of more than $K$ states never win, and
at each switch value exactly 2 sizes tie.
\item $\tau_2=1$, $\tau_3=1+\sqrt2$, $\tau_4=2/(1-2\sqrt2+\sqrt5)=4.906280$ and
$\tau_5=1/(\sqrt3-\frac{1+\sqrt5}2)=8.770636$. The arcs of $2$ states win on an
interval of $\tau$ if and only if $d\ge4$. For $k\ge3$ the arcs of $k$ states
win if and only if $d\ge\lceil3k/2\rceil$.
\item As $k\to\infty$, $\tau_k=\frac{k^3}{2\pi^2}+\frac{3k^2}{4\pi^2}+O(k)$,
and as $d\to\infty$, $\tau_{\rm full}\sim\frac{4d^3}{27\pi^2}$.
\end{enumerate}
\end{theorem}

The cascade stops at two thirds of the cycle. The reason is a comparison per
added state. Growing an arc of $k$ states by one state lowers the exit rate by
$\lambda_k-\lambda_{k+1}\approx2\lambda_k/k$. Closing the cycle removes all of
$\lambda_k$ but adds $d-k$ new states, and this is better once
$\lambda_k/(d-k)>2\lambda_k/k$, that is once $k>2d/3$.

\begin{proof}
By Lemma~\ref{lem:arcs}(a) the admissible faces are the arcs and $C_d$, with
exponents $c_k(\tau)=\tau\lambda_k+k-1$ and $d-1$. Part (a) is
Proposition~\ref{prop:complete}(a) with $d=3$.

(b) Let $d\ge4$. Since $c_k(\tau)-c_{k-1}(\tau)=1-\tau/\tau_k$ and the $\tau_k$
increase, $k\mapsto c_k(\tau)$ is strictly smallest at the $k$ with
$\tau_k<\tau<\tau_{k+1}$, where $\tau_1=0$ and $\tau_d=\infty$. At $\tau_k$
exactly the sizes $k-1$ and $k$ tie. The least arc exponent
$g(\tau)=\min_kc_k(\tau)$ increases continuously and strictly from $0$ to
$\infty$, so $C_d$ alone wins exactly for $\tau$ above the root
$\tau_{\rm full}$ of $g=d-1$. The arcs of $k\ge2$ states
win on an interval if and only if $\tau_k<\tau_{\rm full}$, that is
$c_k(\tau_k)<d-1$, that is $D_k<d$. As the $\tau_k$ increase, the winning sizes
are $1,\dots,K'$ for some $K'$. Let $k_0=\lfloor2d/3\rfloor$, so $2\le k_0\le
d-2$. By Lemma~\ref{lem:arcs}(c), $D_{k_0}<\frac32k_0\le d$ if $k_0\ge3$, and
$D_2=3<4=d$ if $k_0=2$. Also $D_{k_0+1}>\frac32(k_0+1)-\frac12\ge d$, because
$3k_0\ge2d-2$. So $K'=k_0=K$. Then $\tau_K<\tau_{\rm full}<\tau_{K+1}$, so
$c_K(\tau_{\rm full})=d-1$ and $\tau_{\rm full}=(d-K)/\lambda_K=\lceil
d/3\rceil/\lambda_K$. An arc of $k>K$ states has a larger exponent than the
arc of $K$ states for $\tau<\tau_{K+1}$, and one larger than $d-1$ for
$\tau\ge\tau_{\rm full}$, so it never wins. Arcs of 2 sizes tie only at the
distinct values $\tau_k$, and $C_d$ ties with 2 sizes only if $\tau_{\rm full}=\tau_{K+1}$, that is
$D_{K+1}=d$, which Lemma~\ref{lem:arcs}(c) excludes.

(c) The values follow from $\lambda_1=2$, $\lambda_2=1$, $\lambda_3=2-\sqrt2$,
$\lambda_4=(3-\sqrt5)/2$ and $\lambda_5=2-\sqrt3$. By (a) and (b), the arcs of
$k$ states win on an interval if and only if $d>D_k$, and for $k\ge3$ they do
not win at all otherwise. For $k\ge3$ no integer lies in $[D_k,\frac32k)$, so an
integer $d$ exceeds $D_k$ if and only if $d\ge\frac32k$.

(d) Since $2-2\cos x=x^2+O(x^4)$, with $m=k+1$ we get
\[
\lambda_{k-1}-\lambda_k=\pi^2\Bigl(\frac1{(m-1)^2}-\frac1{m^2}\Bigr)+O(m^{-5})
=\frac{2\pi^2}{m^3}\Bigl(1+\frac3{2m}+O(m^{-2})\Bigr).
\]
So $\tau_k=\frac{m^3}{2\pi^2}-\frac{3m^2}{4\pi^2}+O(m)$, which is the
expansion. Also $\lambda_K\sim\pi^2/K^2$ and $K\sim2d/3$, so $\tau_{\rm
full}\sim\frac d3\cdot\frac{4d^2}{9\pi^2}$.
\end{proof}

Table~\ref{tab:cycles} lists the first thresholds. For $d=4,5,6,8$ we get
$\tau_{\rm full}=2$, $2+\sqrt2$, $3+\sqrt5$ and $6+3\sqrt3$.

\begin{table}[t]
\centering\small
\begin{tabular}{@{}crrrr@{}}
\toprule
$k$ & $\lambda_k$ & $\tau_k$ & least $d$ & $b_k$\\
\midrule
2 & 1 & 1 & 4 & $-0.4673558$\\
3 & 0.585786 & 2.414214 & 5 & $-0.9379218$\\
4 & 0.381966 & 4.906280 & 6 & $-1.3514657$\\
5 & 0.267949 & 8.770636 & 8 & $-1.7011216$\\
6 & 0.198062 & 14.308827 & 9 & $-2.0012059$\\
7 & 0.152241 & 21.823898 & 11 & $-2.2632478$\\
8 & 0.120615 & 31.619377 & 12 & $-2.4955119$\\
\bottomrule
\end{tabular}
\caption{Arc thresholds on the cycle $C_d$. The columns $\lambda_k$, $\tau_k$
and $b_k$ do not depend on $d$. The arcs of $k$ states win on an interval of
$\tau$ if and only if $d$ is at least the value in the fourth column
(Theorem~\ref{thm:cycle}(c)), and $b_k$ is the offset of
Corollary~\ref{cor:cycle-second} for the uniform start.}
\label{tab:cycles}
\end{table}


\begin{corollary}[second order on cycles]\label{cor:cycle-second}
Let the start be uniform on $C_d$. Every arc satisfies (H). The arc of $k\le d-1$
states has
\[
K_{\rm arc}(k)=\frac{2}{d(k+1)}\,\frac{(1+\cos\vartheta_k)^2}{\sin^3\vartheta_k}
\Bigl(\frac{k+1}{2\pi}\Bigr)^{(k-1)/2},
\]
and the whole cycle has $K_{C_d}=d^{(d+2)/2}(4\pi)^{-(d-1)/2}$. Let $d\ge4$ and
$2\le k\le\lfloor2d/3\rfloor$, and put
\[
b_k=\tfrac12\log(\lambda_{k-1}-\lambda_k)-\log\frac{K_{\rm arc}(k)}{K_{\rm arc}(k-1)},
\]
which does not depend on $d$. Then for large $N$ the global mode passes from an
arc of $k-1$ states to an arc of $k$ states, and there is $\varepsilon>0$ such
that every change in the size of the arc that carries the global mode with
$\lvert T/L-\tau_k\rvert\le\varepsilon$ happens at $T=\tau_k(L-\frac12\log L+b_k)+o(1)$. Also $b_2=\frac12\log(\pi/8)$.
\end{corollary}

\begin{proof}
Arcs and $C_d$ are connected and have positive start mass, so they satisfy (H),
and Proposition~\ref{prop:constant}(b) applies with $\nu=\mathbf 1$. On an arc of
$k$ states write $\vartheta=\vartheta_k$. The Perron vectors are proportional to
$(\sin j\vartheta)_{j=1}^k$. As $\sum_j\sin^2j\vartheta=(k+1)/2$, we get
$\alpha^*_j=\frac2{k+1}\sin^2j\vartheta$, and we may take
$r=l=\sqrt{\alpha^*}$, so $(\mu\cdot r)(l\cdot\mathbf
1)=\frac1d(\sum_j\sqrt{\alpha^*_j})^2$. The only spanning tree of a path is the
path, so $\tree(c)=\prod_j\alpha^*_j/(\alpha^*_1\alpha^*_k)^{1/2}$. Now
$\sum_j\sin j\vartheta=\cot(\vartheta/2)=(1+\cos\vartheta)/\sin\vartheta$ and
$\alpha^*_1=\alpha^*_k=\frac2{k+1}\sin^2\vartheta$. Also
$\prod_j\sin j\vartheta=(k+1)/2^k$, which follows from
$\prod_{j=1}^k(1-\omega^j)=k+1$ with $\omega=e^{2\pi\mathrm i/(k+1)}$ and
$\lvert1-\omega^j\rvert=2\sin j\vartheta$. Inserting these values gives
$K_{\rm arc}(k)$. On $C_d$ the law $\alpha^*$ is uniform, $r=\mathbf 1$ and
$l=\mathbf 1/d$, so $(\mu\cdot r)(l\cdot\mathbf 1)=1$. Every edge has
conductance $1/d$, and the $d$ spanning trees give $\tree(c)=d^{2-d}$. So
$K_{C_d}=(4\pi)^{-(d-1)/2}d^{(2-d)/2}d^{d}$.

By Theorem~\ref{thm:cycle}(b) the faces tied at $\tau^*=\tau_k$ are the arcs of
$k-1$ and of $k$ states, and all arcs of one size have the same line. As
$\tau_k(\lambda_{k-1}-\lambda_k)=1$, the lines of Theorem~\ref{thm:window}
satisfy
$\ell_k(t)-\ell_{k-1}(t)=\log(K_{\rm arc}(k)/K_{\rm arc}(k-1))+\frac12\log\tau_k+t$,
which vanishes exactly at $t=b_k$. So Theorem~\ref{thm:window} gives the claim.
The factor $1/d$ cancels in $b_k$, and $K_{\rm arc}(1)=1/d$ and $K_{\rm
arc}(2)=4/(d\sqrt{2\pi})$ give $b_2=\log(\sqrt{2\pi}/4)$.
\end{proof}

The first step is the binary problem exactly. By Proposition~\ref{prop:pair},
for every $N$ and $d$ a balanced composition of 2 adjacent states beats a
vertex if and only if $S_N(u)>1$ with $u=h/(1-2h)$, as in
\cite[Thm.~5.1]{paperA}, and $b_2$ is the constant $c_0$ of \cite{paperA}.
Figure~\ref{fig:cycles}(b) shows the support of the exact mode on $C_5$ at
$N=40$ and $72$. The phases come in the order that first order predicts, but
much earlier than $\tau_k\log N$. Most of the gap is the shift $-\frac12\log L+b_k$. At
$N=40$, for example, $\tau_3(L-\frac12\log L+b_3)=5.07$, while the exact switch
from $2$ to $3$ states is at $T=5.12$ and $\tau_3\log N=8.91$.

\begin{remark}[the 4-cycle]\label{rem:C4}
On $C_4$ with the uniform start the winners are a vertex, then an adjacent pair
from $\tau=1$, then $C_4$ from $\tau=2$. Corollary~\ref{cor:cycle-second} gives
$K_{\rm pair}=1/\sqrt{2\pi}$, $K_{\rm arc3}=(4+3\sqrt2)/(4\pi)$ and
$K_{C_4}=8/\pi^{3/2}$, with $\det\Sigma=1/512$ for $C_4$. So by
Theorem~\ref{thm:crossing} every crossing of the pair and $C_4$ satisfies
$T+\log T=2L-\log(8\sqrt2/\pi)+O(1/L)$. The 3-arc never holds the mode at first
order, but Theorem~\ref{thm:crossing} still applies to its face maximum. At
third order, $c=-\frac18$ for the pair by Proposition~\ref{prop:thirdorder}(iii).
For the 3-arc and the whole cycle the formula of
Proposition~\ref{prop:thirdorder}(i) is proved, and its value was evaluated
exactly by computer (Section~\ref{sec:verify}). This gives $c_{\rm
arc3}=\frac54\sqrt2-2$ and $c_{C_4}=-\frac9{16}$. So by
Proposition~\ref{prop:thirdorder}(ii), $T_N$ times the error in the pair-to-cycle
equation tends to $-\frac9{16}+\frac18=-\frac7{16}$.
\end{remark}

\begin{remark}[paths]\label{rem:paths}
The two-thirds rule also holds on the path $P_d$ ($d\ge3$, unit rates). The
winners are the 2 end-arcs of $k$ states for $k=1,\dots,\lfloor2d/3\rfloor$ in
turn, and then $P_d$. Reflecting an end-arc of $k$ states at its leaf gives a
path of $2k$ states, whose Perron vector is symmetric. So the end-arc has exit
rate $\lambda'_k=\lambda_{2k}=2-2\cos(\pi/(2k+1))$, below the rate $\lambda_k$
of an inner arc, and it is strictly decreasing and convex in $k$ by
Lemma~\ref{lem:arcs}(b). The proof of Theorem~\ref{thm:cycle} then goes through
with $\lambda'_k$, because
$D'_k=k+\lambda'_k/(\lambda'_{k-1}-\lambda'_k)$ satisfies
$\frac32k-\frac12<D'_k<\frac32k$ for all $k\ge2$. As in
Lemma~\ref{lem:arcs}(c), the lower bound follows from
$(2k-1)^2(k+1)-(k-1)(2k+1)^2=2$. For the upper bound put
$y=(8k-2)/((k+2)(2k-1)^2)\le3$, so that $\frac
k{k+2}=(1+y)(\frac{2k-1}{2k+1})^2$. Using $\sqrt{1+y}\ge1+y/3$, the sinc bounds
and $1-\pi^2/216>0.954$, it reduces to
$\frac{\pi^2}{12k}+\frac{\pi^4}{7680k^2}<0.318\,\frac{8k-2}{k+2}$, which holds at
$k=2$ ($0.414<1.113$) and hence for all $k\ge2$, since the left side decreases
and the right side increases. The first switch is at the
golden ratio $1/(\lambda'_1-\lambda'_2)=(1+\sqrt5)/2$.

The lower bound is nearly attained, and this is visible at finite $N$. On $P_4$,
$D'_3=4.077>4$, so the end-arc of $3$ states never wins at first order, but in
the plot of Lemma~\ref{lem:symmetric}(c) its point $(2,\underline\lambda_3)$
lies only $0.0071$ above the lower hull. Numerically, the exact law on $P_4$
has this end-arc as its mode for $T$ in $(9.381,12.502)$, $(12.176,15.142)$,
$(15.804,18.443)$ and $(19.292,21.590)$ at $N=40$, $72$, $150$ and $300$, a
window that shrinks only slowly. We also computed $C_4$, $C_5$, $C_6$ and the
paw of Remark~\ref{rem:nonnested} exactly at some $N$ between $40$ and $300$.
On these graphs every size that never wins has its point
$(s-1,\underline\lambda_s)$ at least $0.077$ above the lower hull (on the paw
every size wins), and the exact phases came in the order that first order
predicts. So at feasible $N$ the order of the phases can differ from first
order when such a point lies close to the hull.
\end{remark}

\subsection{Two triangles}\label{sec:ex-triangles}

Let $\Theta_\varepsilon$ have states $1,\dots,6$, unit rates on the triangles
$A=\{1,2,3\}$ and $B=\{4,5,6\}$, and one more link $3$--$4$ of rate
$\varepsilon>0$. We call $1,2,5,6$ the inner states and put
$p_A(x)=x^2-(3+\varepsilon)x+\varepsilon$. For small $\varepsilon$ each triangle
is a well-separated community. It holds the mode from $\tau\approx1$ until
$\tau\approx9/\varepsilon$ (Proposition~\ref{prop:triangles}).

\begin{proposition}[two triangles joined by a weak link]\label{prop:triangles}
Let the start have full support.
\begin{enumerate}[(a)]
\item $\lambda_A=\lambda_B=\frac12\bigl(3+\varepsilon-\sqrt{\varepsilon^2+2\varepsilon+9}\bigr)
=\frac\varepsilon3-\frac{2\varepsilon^2}{27}+O(\varepsilon^3)$ is the smaller
root of $p_A$, and $0<\lambda_A<1$.
\item The least exit rates by size are $\underline\lambda_1=2$ (the inner states), $\underline\lambda_2=1$
(the inner edges $\{1,2\}$ and $\{5,6\}$), $\underline\lambda_3=\lambda_A$ (the triangles),
$\underline\lambda_4=\lambda_{\{1,2,3,4\}}$ (and its mirror image $\{3,4,5,6\}$),
$\underline\lambda_5=\lambda_{\{2,\dots,6\}}$ (and its images under the symmetries), and
$\underline\lambda_6=0$. Here $\underline\lambda_4$ and $\underline\lambda_5$ are the least roots of
$p_4(x)=x^3-(5+2\varepsilon)x^2+(6+6\varepsilon)x-2\varepsilon$ and
$p_5(x)=x^3-(4+2\varepsilon)x^2+(3+4\varepsilon)x-\varepsilon$.
\item Let $\varepsilon_1=(64+\sqrt{46})/60=1.17971$, let
$\varepsilon_2=4.48108$ be the positive root of
$\varepsilon^3+2\varepsilon^2-19\varepsilon-45$, and put
\[
\tau_A=\frac1{1-\lambda_A}=\frac{1+\varepsilon+\sqrt{\varepsilon^2+2\varepsilon+9}}4,
\qquad
\tau^*=\frac3{\lambda_A}=\frac3{2\varepsilon}\Bigl(3+\varepsilon+\sqrt{\varepsilon^2+2\varepsilon+9}\Bigr).
\]
For $\varepsilon<\varepsilon_1$ the winners are the inner states, then the inner
edges from $\tau=1$, then the triangles from $\tau_A$, and then $[6]$ from
$\tau^*$. For $\varepsilon_1<\varepsilon<\varepsilon_2$ a phase of
$\{1,2,3,4\}$ and its mirror image comes between the triangles and $[6]$, with
switches at $1/(\lambda_A-\underline\lambda_4)$ and $2/\underline\lambda_4$. For
$\varepsilon>\varepsilon_2$ the triangles never win, and the inner edges pass to
$\{1,2,3,4\}$ and its mirror image at $2/(1-\underline\lambda_4)$ and then to $[6]$ at
$2/\underline\lambda_4$. As $\varepsilon\to0$,
$\tau_A=1+\frac\varepsilon3+O(\varepsilon^2)$ and
$\tau^*=\frac9\varepsilon+2+\frac{2\varepsilon}9+O(\varepsilon^2)$.
\end{enumerate}
\end{proposition}

\begin{proof}
(a) On vectors $(a,a,b)$ the matrix $-Q_A$ acts as
$(a,b)\mapsto(a-b,\,-2a+(2+\varepsilon)b)$, with characteristic polynomial
$p_A$, and $(1,-1,0)$ has eigenvalue $3$. As $p_A(0)=\varepsilon>0$ and
$p_A(1)=-2<0$, the smaller root of $p_A$ lies in $(0,1)$, and it is $\lambda_A$.
The expansion follows from
$\sqrt{9+2\varepsilon+\varepsilon^2}=3+\frac\varepsilon3+\frac{4\varepsilon^2}{27}+O(\varepsilon^3)$.

(b) The admissible faces are the connected ones. For a face $F$, $-Q_F$ is the
Laplacian of the graph on $F$, which is positive semidefinite, plus the diagonal
matrix of the rates out of $F$. So $\lambda_F\ge1$ whenever every state of $F$
has rate at least $1$ out of $F$. This covers $\{3,4\}$, the pairs $\{1,3\}$,
$\{2,3\}$ and their mirror images (for which in fact
$\lambda=\frac12(4+\varepsilon-\sqrt{\varepsilon^2+4})>1$), every connected
triple other than $A$ and $B$, and the sets $\{a,3,4,b\}$ with $a\in\{1,2\}$
and $b\in\{5,6\}$. The other connected faces are the inner states (exit rate
$2$, against $2+\varepsilon$ for $3$ and $4$), the inner edges ($\lambda=1$),
$A$ and $B$, $\{1,2,3,4\}$ and its mirror image, the complements of the inner
states, which are all images of $\{2,\dots,6\}$, and $[6]$. Since
$\{2,\dots,6\}$ contains the mirror image $\{3,4,5,6\}$ of $\{1,2,3,4\}$,
Lemma~\ref{lem:symmetric}(b) gives
$\underline\lambda_5<\lambda_{\{1,2,3,4\}}<\lambda_A<1$. Comparing within each
size gives (b), once we know the polynomials. An exact computer calculation
(Section~\ref{sec:verify}) gives $\det(xI+Q_F)=(x-3)p_4(x)$ for
$F=\{1,2,3,4\}$ and $(x-3)^2p_5(x)$ for $F=\{2,\dots,6\}$. These are
determinants of size $4$ and $5$, and they can also be checked by hand. So
$\underline\lambda_4$ and $\underline\lambda_5$ are the least roots of $p_4$
and $p_5$.

(c) Put $w_s=(s-1,\underline\lambda_s)$. Suppose that a convex piecewise
linear function through some of the points $w_s$ lies strictly below the
others. Then by Theorem~\ref{thm:hull}(a) and Lemma~\ref{lem:symmetric}(c) the
sizes at its corners win in turn. The switch between 2 consecutive corners is
at $-1/\sigma$, where $\sigma$ is the slope between them. We need 3
comparisons, namely (i) $\underline\lambda_5>\underline\lambda_4/2$ for every
$\varepsilon>0$, (ii) $\underline\lambda_4>\frac23\lambda_A$ if and only if
$\varepsilon<\varepsilon_1$, and (iii)
$\lambda_A<\frac12(1+\underline\lambda_4)$ if and only if
$\varepsilon<\varepsilon_2$. The 3 differences are continuous in
$\varepsilon$. If one of them vanishes, then $\underline\lambda_5$ is a common
root of $p_5(x)$ and $p_4(2x)$, or $\underline\lambda_4$ is a common root of
$p_4(x)$ and $p_A(\frac32x)$, or $\underline\lambda_4$ is a common root of
$p_4(x)$ and $p_A(\frac{1+x}2)$. An exact computer calculation
(Section~\ref{sec:verify}) gives the resultants
\begin{align*}
\operatorname{Res}_x\bigl(p_5(x),p_4(2x)\bigr)&=-8\varepsilon^3(16\varepsilon^2+12\varepsilon+9),\\
\operatorname{Res}_x\bigl(p_A(\tfrac32x),p_4(x)\bigr)&=-\tfrac1{16}\varepsilon^2(120\varepsilon^2-256\varepsilon+135),\\
\operatorname{Res}_x\bigl(p_A(\tfrac{1+x}2),p_4(x)\bigr)&=\tfrac18(\varepsilon^3+2\varepsilon^2-19\varepsilon-45).
\end{align*}
So the difference in (i) never vanishes, the one in (ii) can vanish only at
$\varepsilon_0=(64-\sqrt{46})/60=0.95363$ and at $\varepsilon_1$, and the one in
(iii) only at $\varepsilon_2$, by Descartes' rule of signs.

We find the signs near $0$ and $\infty$. As $\varepsilon\to0$,
$\lambda_A\to0$ and so $\underline\lambda_4,\underline\lambda_5\to0$. At $\varepsilon=0$ the cubics are
$x(x-2)(x-3)$ and $x(x-1)(x-3)$, so the implicit function theorem at the simple
root $0$ gives $\underline\lambda_4=\frac\varepsilon3+O(\varepsilon^2)$ and
$\underline\lambda_5=\frac\varepsilon3+O(\varepsilon^2)$. With (a) the differences in (i),
(ii) and (iii) are $\frac\varepsilon6+O(\varepsilon^2)$,
$\frac\varepsilon9+O(\varepsilon^2)$ and $\frac12+o(1)$, all positive. As
$\varepsilon\to\infty$, $\lambda_A\to1$, while
$\underline\lambda_4\le(3-\sqrt5)/2<\frac12$ for every $\varepsilon$. This bound comes from
Lemma~\ref{lem:symmetric}(a) with the test vectors $(a,a,c,c)$ on
$\{1,2,3,4\}$, whose quotient $((a-c)^2+c^2)/(a^2+c^2)$ does not involve
$\varepsilon$ and has least value $(3-\sqrt5)/2$. So the differences in (ii) and
(iii) are negative for large $\varepsilon$. This proves (i) and (iii). It also
proves (ii) once we know that the difference in (ii) is nonzero at
$\varepsilon_0$, since then $\varepsilon_1$ is its only possible zero. For that
put $x_0=7/40$. Then
\[
p_4(x_0)=\frac{57743}{64000}-\frac{809}{800}\varepsilon,\qquad
p_A(\tfrac32x_0)=\frac{59}{80}\varepsilon-\frac{4599}{6400}
\]
are both negative for $0.8923<\varepsilon<0.9743$, an interval that contains
$\varepsilon_0$. On $[0,\frac12]$ we have $p_4'(x)\ge1+4\varepsilon>0$, and
$p_4(0)<0$, so $p_4$ has no root in $[0,x_0]$ and $\underline\lambda_4>x_0$. Also
$p_A(\frac32x_0)<0$ puts $\frac32x_0$ between the roots of $p_A$, so
$\lambda_A<\frac32x_0$. Hence $\underline\lambda_4>\frac23\lambda_A$ at $\varepsilon_0$.

Now we read off the hull. As $p_A(\frac34)=\frac\varepsilon4-\frac{27}{16}$, we
have $\lambda_A<\frac34$ for $\varepsilon<\varepsilon_1<\frac{27}4$. For
$\varepsilon<\varepsilon_1$ the corners are $w_1,w_2,w_3,w_6$, with slopes $-1$,
$\lambda_A-1$ and $-\lambda_A/3$, which increase because $0<\lambda_A<\frac34$.
The segment $w_3w_6$ has height $\frac23\lambda_A$ at $x=3$ and
$\frac13\lambda_A$ at $x=4$, so $w_4$ and $w_5$ lie above it by (i) and (ii).
For $\varepsilon_1<\varepsilon<\varepsilon_2$ the corners are
$w_1,w_2,w_3,w_4,w_6$, with slopes $-1$, $\lambda_A-1$, $\underline\lambda_4-\lambda_A$ and
$-\underline\lambda_4/2$. These increase because $\lambda_A>0$, then by (iii), then by (ii),
and $w_5$ lies above $w_4w_6$ by (i). For $\varepsilon>\varepsilon_2$ the corners
are $w_1,w_2,w_4,w_6$, with slopes $-1$, $(\underline\lambda_4-1)/2$ and $-\underline\lambda_4/2$, which
increase because $\underline\lambda_4<\frac12$. Here $w_3$ lies above $w_2w_4$ by (iii), and
$w_5$ above $w_4w_6$ by (i). The switch values are minus the reciprocals of the
slopes. The closed forms of $\tau_A$ and $\tau^*$ follow by rationalizing, and
their expansions from (a).
\end{proof}

So for small $\varepsilon$ the inner edges win only on the short interval
$(1,\tau_A)$, and a triangle then keeps the mode until $\tau^*$. For
$\varepsilon>\varepsilon_1$ the link is no longer weak enough, and the mode takes
in the far end of the link before it spreads.

\begin{remark}[the mode can leave a community]\label{rem:nonnested}
The winners need not grow around the first vertex. Take the paw, a triangle
$\{1,2,3\}$ with a leaf $4$ joined to $1$, with unit rates and a start of full
support. The single states have exit rates $3$ for $\{1\}$, $2$ for $\{2\}$
and $\{3\}$, and $1$ for $\{4\}$. The connected pairs have $2-\sqrt2$ for
$\{1,4\}$, $1$ for $\{2,3\}$ and $(5-\sqrt5)/2$ for $\{1,2\}$ and $\{1,3\}$.
The triangle $\{1,2,3\}$ has $2-\sqrt3$, from the vectors $(c,e,e)$ with
characteristic polynomial $x^2-4x+1$. The triples $\{1,2,4\}$ and
$\{1,3,4\}$ have $2-2\cos(2\pi/9)=0.468$, the least root of
$x^3-6x^2+9x-3$. So
$(\underline\lambda_1,\dots,\underline\lambda_4)=(1,2-\sqrt2,2-\sqrt3,0)$. The
slopes $1-\sqrt2<\sqrt2-\sqrt3<\sqrt3-2$ increase, since $(1+\sqrt3)^2<8$ and
$(2+\sqrt2)^2<12$. So by Theorem~\ref{thm:hull}(a) every size wins, in the
order $\{4\}$, then $\{1,4\}$ from $\tau=1+\sqrt2$, then $\{1,2,3\}$ from
$\sqrt2+\sqrt3$, then all $4$ states from $2+\sqrt3$. The mode leaves the
leaf, since $\{1,4\}\not\subseteq\{1,2,3\}$.
\end{remark}
